\section{Summary}

This thesis introduced \Natural, a sovereign and portable NL2SQL pipeline using
multiple techniques such as schema-aware example selection ($\sigma$), schema
subsetting ($\phi$), query projection ($\pi$), self-refinement ($\rho$) and
consensus voting ($\nu$) while remaining executable on consumer-grade systems.
It was evaluated against multiple variants of \OmniSQL~and GPT-4 on \Spider~and
\Bird~and a brief ablation study was performed to identify the impact of
pipeline components and ICL example sources. The results are presented in the
evaluation and discussion sections respectively. The key contributions of this
thesis are:

\begin{enumerate}
    \item \textbf{Graph-Based Example Selection} — A novel schema-aware example
        selection mechanism using Wasserstein-Weisfeiler-Lehman graph kernels
        to measure the structural similarity between two database schemas.
    \item \textbf{Portable Pipeline Design} — A five-stage pipeline design
        combining example selection ($\sigma$), schema subsetting ($\phi$),
        few-shot generation ($\pi$), self-refinement ($\rho$), and consensus
        voting ($\nu$) that demonstrated substantial accuracy improvements in
        a local, controlled environment.
    \item \textbf{Production Implementation} — A complete, production-grade Rust
        implementation optimized for consumer-grade hardware (RTX 3090, 24GB
        VRAM), achieving a latency of 7.3-16.0s on \Spider~and 9.3-43.8s on
        \Bird~depending on the configuration used.
    \item \textbf{Highlighted Reproducibility Gap} — Documentation and
        discussion of the 28.1pp performance gap between published and
        reproduced results of \OmniSQL, highlighting critical challenges in
        LLM-based research reproducibility and transferability of models,
        prompting future research.
\end{enumerate}

\subsection{Research Objectives and Achievement}

The primary research objective of this thesis was to design and implement a
production-ready, portable and sovereign NL2SQL system which avoids the
dependency on proprietary APIs for inference. This goal was achieved by
introducing \Natural, which showed that small and portable LLMs, paired with
schema-aware example selection can achieve meaningful performance on imperfect,
real-world databases. In local evaluations \Natural~performed +2.0-15.8pp
better than the \OmniSQL-7B-gguf model, but a gap remains to the published
\OmniSQL~results. This is subject to further evaluations as a full
reproduction of the reported baseline model performance was not possible and a
significant gap of up to 28.1pp was encountered. Using techniques such as
ICL combined with schema-aware example selection can partially close the
performance gap to proprietary solutions while maintaining full sovereignty and
portability.

\subsection{Research Questions}

The research questions formulated in section~\ref{section:introduction:research-questions}
have been addressed in the discussion section. The systematic
evaluation gathered and analyzed the performance characteristics and indicates
that open-source models can achieve competitive NL2SQL performance. In the
locally controlled environment, systematic pipeline composition with
graph-based example selection showed to be the critical performance driver,
especially on complex databases. Pipeline components show synergistic rather
than additive tendencies, but further evaluations, especially better ablations
are required to make an absolute assessment.

Notable tradeoff between accuracy and latency was observed across different
configurations of \Natural, which informs deployment decisions and viability
for different use cases.

\subsection{Future Work}

The work on \Natural~revealed several important areas that subsequent research
should focus on. These areas include system improvements, reproducibility
research, and user-oriented research.

\subsubsection{Reproducibility Research}

Based on the performance gap analysis in section~\ref{sec:eval:performance-gap},
several recommendations for further research in this area can be made:

\paragraph{Systematic Quantization Study}

A dedicated investigation of the impact of model quantization on NL2SQL
performance could provide guidance for better understanding the quantization-to-performance
tradeoffs. This would allow for empirical decision-making when
choosing model sizes and quantization for resource-constrained
environments. Prior research exists on the general impact of quantization on
model performance, but having concrete datapoints specific to NL2SQL would
clarify the observed behavior \citep{Int8AtScale}. This study should evaluate
SOTA NL2SQL systems across multiple quantization levels (F16, Q8,
Q6, Q4, and Q3) on all major NL2SQL benchmarks such as \Spider, \Spider~2.0, and
\Bird~measuring both the accuracy and latency implications of different
configurations. Furthermore, an analysis could reveal correlation in performance
behavior between quantization levels and query and schema complexity.

\paragraph{Cross-Framework Validation}

To isolate inference engine effects and subtle differences introduced by
different model formats, the same model should be evaluated across multiple
frameworks. This study could include PyTorch and Transformers, llama.cpp, vLLM,
and ONNX and evaluate SOTA NL2SQL models in varying sizes against
prevalent benchmarks. Consistent results across frameworks would further
increase confidence in reproducibility of results, while discrepancies would
highlight implementation-specific details of algorithms and
prompting.

\paragraph{Schema and Example Presentation}

While research on schema and example presentation exists within the NL2SQL
community, it is not exhaustive and includes evaluations of different classes
of schema and query presentation \citep{DAIL-SQL, XiYan}. Further research in
this field could highlight differences introduced through subtle changes in
presentation of examples or schemas and recommend mitigation techniques such as
identifying the most performant presentation formats or fine-tuning base models
to a specific presentation format.

\subsubsection{System Improvements}

Significant opportunities for future research revolve around further system and
performance improvements to enable deployments on less powerful hardware. The
overall candidate latency of \Natural~emerged as a critical blocker for
real-world adoption, as requiring users to procure expensive research or
enterprise-grade GPUs is suboptimal. Substantial performance optimization is
needed to achieve a democratization effect.

\paragraph{Model Distillation}

Through further research on NL2SQL models and the potential of model
distillation, significantly smaller LLMs could yield a noticeable increase in
performance and drop the hardware requirements. For users with existing
high-end hardware, this could yield a significant candidate latency reduction
and unlock interactive use cases that are not viable with the latency values
measured in this thesis.

\paragraph{Component Refinement}

Additional component refinement and removing components
that add minimal or no additional performance would both increase
the semantic accuracy of \Natural~and reduce potentially needless
computational overheads. Especially the schema subsetting component has shown
little to no notable performance improvement, and could potentially even
harm performance.

\paragraph{Multi-Model Architectures}

Lastly, this thesis did not utilize a multi-model or agentic architecture as
proposed by \cite{XiYan}. Future research could integrate the learnings of this
thesis with a multi-model or agentic architecture to see whether a
reimplementation of \Natural's components as tool calls provides a
performance gain over agentic architectures without example selection
algorithms and self-refinement.

\subsection{Closing Remarks}

This thesis demonstrates that portable, sovereign NL2SQL systems based on
open-source models represent a promising alternative to systems relying on
proprietary APIs, making NL2SQL deployments in constrained domains such as
government, healthcare, and finance feasible. Significant challenges remain,
particularly regarding latency, semantic accuracy on complex databases, and
reproducibility of research results. The algorithms introduced in this thesis,
especially the schema-aware example selection, prove that additional
algorithmic support systems can partially close the performance gap between
open-source and closed-source models in controlled environments. Contributions
were made both through a practical implementation of \Natural and
methodological insights for future NL2SQL research, particularly by
highlighting the importance of example sources and selection algorithms. As
foundational models continue to improve and hardware becomes more accessible,
NL2SQL systems will play an increasingly important role in democratizing data
access while maintaining data sovereignty.
